\section{Trò chơi dựng trên attribution}
\label{sec:ch11-tro-choi-dung-tren-attribution}

\index{metagame|(}%
Hai nhóm phương pháp có sẵn để lại một yêu cầu kép: tách sạch
\tn{hiệu ứng riêng}{pure individual effect} như nhóm theo tập làm được, mà vẫn
giữ được chiều như nhóm nối tầng có. Xây dựng của bài báo đáp ứng cả hai, và nó
làm điều đó bằng một phép chuyển mà chương~\ref{ch:shap} đã dựng sẵn công cụ.

Phép chuyển ấy là như sau. Mục~\ref{sec:ch04-tro-choi-chia-gia-tri} đã trình
bày giá trị Shapley như lời giải của bài toán chia phần thưởng trong một trò
chơi hợp tác: có một tập người chơi, có một hàm gán cho mỗi liên minh một giá
trị, và giá trị Shapley chia giá trị của liên minh đầy đủ cho từng người chơi.
Khi SHAP dùng công cụ ấy để giải thích một mô hình, người chơi là các đặc trưng
còn hàm giá trị là đầu ra của mô hình trên từng liên minh. Bài báo này giữ
nguyên công cụ và thay hàm giá trị: hàm giá trị không còn là mô hình nữa mà là
chính phương pháp attribution. Trò chơi mới ấy là cái mà các tác giả gọi là
metagame.

Nói cho đủ thì phải nói cả một bước nữa, và bước ấy là chỗ khác nhau giữa xây
dựng này với phép nối tầng của phương trình~\ref{eq:ch11-noi-tang}. Muốn đo ảnh
hưởng của đặc trưng $j$ lên attribution của đặc trưng $i$, ta giữ nguyên giá
trị của $i$ ở mọi liên minh thay vì để nó thay đổi cùng $j$. Phép nối tầng
không làm vậy: ở đó cả $x_i$ lẫn $x_j$ cùng biến thiên, nên hiệu ứng riêng và
hiệu ứng cặp dính vào nhau, và bài chỉ ra đó chính là nguồn của chỗ rò mà
bảng~\ref{tab:ch11-tach-hieu-ung} đo được~\autocite{p23metagame}. Cố định $x_i$
rồi mới lấy giá trị Shapley của $j$ trên $d-1$ đặc trưng còn lại là điều làm
cho con số thu được có chiều.

Định nghĩa dưới đây cần thêm một ký hiệu. Viết $\varphi_i(S; x, f)$ cho
attribution của đặc trưng $i$ khi chỉ những đặc trưng trong tập $S$ có mặt, các
đặc trưng ngoài $S$ bị thay bằng baseline. Đó là cách che đặc trưng mà
mục~\ref{sec:ch08-vong-lap-chung} đã đặt tên và là quy ước duy nhất mà bài
dùng; $\varphi_i(x,f)$ trong sách là trường hợp $S$ chứa tất cả.

\begin{definition}[Meta-attribution]
  \label{def:ch11-meta-attribution}
  Cho một phương pháp attribution bậc nhất $\varphi_i(x,f)$ và hai đặc trưng
  $i \neq j$. Meta-attribution $\varphi_{j \to i}(x,f)$, tức ảnh hưởng của đặc
  trưng $j$ lên attribution của đặc trưng $i$, được định nghĩa là
  %
  \begin{equation*}
    \varphi_{j \to i}(x,f) :=
    \sum_{S \subseteq [d] \setminus \{i,j\}}
      \frac{1}{(d-1)\binom{d-2}{|S|}}
      \Bigl[\varphi_i\bigl(S \cup \{i,j\}; x, f\bigr)
          - \varphi_i\bigl(S \cup \{i\}; x, f\bigr)\Bigr].
  \end{equation*}
  %
  Hiệu ứng riêng là giá trị biên của cùng công thức ấy,
  $\varphi_{i \to i}(x,f) := \varphi_i(\{i\}; x, f)$~\autocite{p23metagame}.
\end{definition}

Từng phần của định nghĩa~\ref{def:ch11-meta-attribution} đọc được bằng những gì
Phần~II đã dựng. Trong ngoặc vuông là đóng góp biên của $j$: attribution mà
đặc trưng $i$ nhận được khi liên minh có cả $j$, trừ đi attribution mà nó nhận
được khi liên minh ấy không có $j$. Cả hai lần, $i$ đều nằm trong liên minh,
nên đại lượng bị trừ luôn là attribution của cùng một đặc trưng $i$ trên hai
tình huống chỉ khác nhau ở sự có mặt của $j$. Trọng số phía trước là trọng số
Shapley của một trò chơi $d-1$ người chơi, vì $i$ đã bị giữ cố định và không
còn tham gia. Nó là cùng một trọng số mà
phương trình~\ref{eq:ch04-gia-tri-shapley} viết bằng giai thừa, chỉ chép sang
dạng nhị thức: với $n$ người chơi, $|S|!\,(n-|S|-1)!\,/\,n!$ bằng
$1/\bigl(n\binom{n-1}{|S|}\bigr)$, và ở đây $n = d-1$. Tổng chạy trên mọi liên
minh $S$ không chứa $i$ và $j$.

Ký hiệu $\varphi_{j \to i}$ nối vào ký hiệu $\varphi_j$ mà
phụ lục~\ref{app:ky-hieu} đã có, và nối một cách có chủ ý: cái được lấy đóng
góp biên ở đây chính là attribution của định nghĩa~\ref{def:ch06-attribution},
nên đại lượng mới là một mở rộng của đại lượng cũ chứ không phải một đại lượng
trùng tên. Mũi tên trong chỉ số đọc theo chiều ảnh hưởng, từ đặc trưng gây ảnh
hưởng sang attribution chịu ảnh hưởng.

Khi phương pháp bậc nhất là giá trị Shapley thì bài gọi kết quả là Meta-Shapley
value, viết tắt Meta-SV; khi nó là Integrated Gradients hay gradient nhân đầu
vào thì tên là Meta-IG và Meta-G$\times$I. Ba tên ấy là ba dòng cuối của
bảng~\ref{tab:ch11-tach-hieu-ung}, và cả ba đều trả đúng $x_1$ cùng $0$ cho hai
hiệu ứng riêng, đồng thời cho hai giá trị khác nhau ở hai chiều của cặp. Hai
trong ba dòng ấy còn khác nhau ở chỗ đáng chú ý: Meta-SV cho
$\varphi_{2\to1} = \varphi_{1\to2} = \tfrac{1}{2}T$, tức đối xứng, còn Meta-IG
cho $\tfrac{1}{3}T$ và $\tfrac{2}{3}T$, tức không đối xứng. Bài giải thích rằng
tính đối xứng hay không là tính chất của phương pháp bậc nhất nằm dưới chứ
không phải của xây dựng bậc hai, và dù đối xứng hay không thì cả ba biến thể
đều tách được hiệu ứng riêng khỏi hiệu ứng tương tác~\autocite{p23metagame}.
\index{metagame|)}%
